\chapter{Evaluation}
\label{ch:eval}

This chapter reports what the reference implementation actually demonstrates.
We are careful to separate verified behavior from scaffolding.

\section{FFI$\leftrightarrow$Lean Consistency}
\label{sec:eval-consistency}

The theorem calculator (\cref{ch:ffi}) verifies, for each pattern at $n=10$,
that the C{--} kernel and Lean agree. In a passing run the gate reports:
\begin{verbatim}
Verifying FFI <-> Lean consistency...
  [1] SumSquares(10): consistent
  [2] SumLinear(10): consistent
  [3] SumCubes(10): consistent
=== All verifications passed ===
\end{verbatim}
This is a real, executable cross-check: two independent implementations of
three closed forms agree on tested inputs. It is not a proof of the closed
forms (that is Lean's job at Layer~4), but it is strong evidence that the
kernel is not silently wrong.

\section{Benchmark: SumSquares(1..1000)}
\label{sec:eval-bench}

The driver benchmarks 1000 FFI computations of SumSquares: it reports total
milliseconds and the final result. On the reference hardware this exercises
the arena allocator and rewrite dispatch 1000 times consecutively. While
absolute numbers depend on the build, the \emph{existence} of a reproducible
benchmark in the verified path is the point: performance is measured on the
trusted path, not a separate fast-and-unchecked one.

\section{Sample Results}
\label{sec:eval-samples}

For $n \in \{1,2,5,10,50,100\}$ the calculator prints
$\ensuremath{Sigma} k^2$, $\ensuremath{Sigma} k$, $\ensuremath{Sigma} k^3$ from the kernel. These agree with the
analytic formulas by construction of the dispatch in \texttt{bridge.c}; the
Lean side independently re-derives them via \texttt{patternRegistry}, so the
printed table is a visible instance of the cross-validation.

\section{FOL Theorem Prover}
\label{sec:eval-fol}

The Layer~6 FOL checker reports \texttt{15/15 theorems verified}. This is an
internal consistency result: a resolution prover confirms fifteen
first-order statements. It complements Lean and the kernel as a third
independent line of verification.

\section{What Is Not Yet Evaluated}
\label{sec:eval-gaps}

Honesty requires listing what the current system does \emph{not} yet measure:
\begin{itemize}
  \item End-to-end latency through all ten layers (the orchestration glue is
        scaffolding, \cref{ch:build}).
  \item Completeness score $C$ across the full theorem corpus (only three
        closed forms are currently formalized).
  \item Large-scale swarm throughput in \texttt{agentos}.
\end{itemize}
These are active frontiers, tracked by \texttt{sorrydb} and the build gate.

\section{Reported Results}
\label{sec:eval-results}

The following table summarizes the measurements available on the
reference target used for this paper (single-core x86-64, in-tree
toolchain).

\begin{table}[ht]
\centering
\caption{Measured pipeline characteristics (reference runs).}\label{tab:results}
\small
\begin{tabular}{l r r}
\toprule
\textbf{Metric} & \textbf{Value} & \textbf{Source} \\
\midrule
Layers in chain & 10 & \cref{ch:vscp} \\
Closed forms certified & 3 & \texttt{TheoremCalc} \\
Build-gate halts on failure & yes & \cref{ch:build} \\
Receipt hash algorithm & SHA-256 & \cref{ch:worm} \\
Plasma Gate signature & Ed25519 & \cref{ch:security} \\
Bitcoin anchor prefix & \texttt{MATH5:} & \cref{ch:anchors} \\
Repository stars & 3 & GitHub \\
\bottomrule
\end{tabular}
\end{table}

These are conservative, single-run figures; the architecture is designed so
that the \emph{correctness} claims hold independently of the numeric values,
which are presented only to characterize the implementation.
